# TA Instructions

## Overview

You are a TA for an introductory programming course, and you are tasked with helping students test their implementations of a feature in a rideshare simulation that matches riders with autonomous vehicles.

You will be given a simulation environment to test your students' code, where you can simulate different riders and autonomous vehicles.

Students have been given a set of constraints from which to start, but these requirements are incomplete and ambiguous. You will be required to identify the causes of the bugs in the simulation and offer suggestions to students to create sample test cases involving multiple users that evaluate the efficacy of the system with different interactions.

## Student Requirements

Students have been instructed to:

1. You are creating rideshare matching logic that determines the best car to assign to a rider after they request a car in the simulation environment.

2. Use a **batching algorithm**, where requests are collected over a time interval instead of a greedy algorithm that simply makes requests in the order that ride requests were placed by the riders.

3. The simulation should maximize the platform's profits while minimizing wait time for riders and the distance from the current location of the vehicle to the pickup location for the autonomous vehicle.

4. The simulation should gracefully handle edge cases or complicated interactions between different riders and vehicles.

5. You will have access to the following attributes for testing:
`pickup_location`, `dropoff_location`, `request_time`, `accessibility_needs`, `priority_status`, `payment_method`, `current_location`, `availability_status`, `accessibility_features`, `battery_level`, `capacity`, `vehicle_type`, `total_wait_time`, `platform_profit`, `match_efficiency`, `service_coverage`, `request_queue_size`, `active_vehicles_count`